\section{The ordinary target}
\label{sec:ordinary}

Throughout this section, let $S=\langle m=g_1<\cdots<g_e\rangle$
be minimally generated, with $e\ge3$, and put
\[
 G=G_S,\qquad \delta=m-e,\qquad B=m-1,\qquad
 \mathcal V=e+\frac\delta2-\frac32.
\]
The number $\mathcal V$ is a resource bound, not the width of $S$.
We prove the following assertion.

\begin{proposition}\label{prop:ordinary}
If $g_e-m\ge m-1$, then, over every field $k$,
\[
 \beta_i(G_S;k)\le\beta_i(G_{\widetilde S};k)
   =i\binom m{i+1}\quad(i\ge1).
\]
The zeroth Betti numbers on both sides are one.
\end{proposition}

\subsection{Coordinates, one cut, and one resource sequence}

The degree-$n$ monomial basis of $G$ consists
of $[v]$ with $\operatorname{ord}(v)=n$; multiplication survives precisely
when orders add.  Set
\[
 x=[m],\quad A=G/xG,\quad T=(0:_Gx),\quad
 S_0=k[Y_2,\ldots,Y_e].
\]
A basis element of $A$ is a \emph{head}: its coordinate has no surviving
incoming $x$-arrow.  Non-Ap\'ery heads span an $S_0$-ideal $N$ in $A$.
The quotient $Q=A/N$ has basis $\operatorname{Ap}(S,m)$ and is standard
graded.  For $q_s=\dim_kQ_s$, $a=q_2$, and
$\nu=\max\{s:q_s\ne0\}$, we have
\begin{equation}\label{eq:ordinary-apery}
 q_0=1,\quad q_1=e-1,\quad \sum_{s\ge2}q_s=\delta,
 \qquad A_2=Q_2,\quad \nu\le\delta+1.
\end{equation}
Indeed, Ap\'ery elements have maximal expressions without $m$, and every
prefix of such an expression realizes its order.  This proves standard
generation and the displayed counts; absence of gaps in the Hilbert
function gives the last inequality.

The monomial basis of $T$ consists of \emph{terminals} $\tau$ whose
outgoing $x$-arrow dies.  Their jumps
$d(\tau)=\operatorname{ord}(\tau+m)-\operatorname{ord}(\tau)$ satisfy
$2\le d(\tau)\le\nu$.  Translation $\tau\mapsto\tau+m$ bijects terminals
with non-Ap\'ery heads.  The appropriate module comparison uses the
opposite filtrations
\[
 D_sT=\langle\tau:d(\tau)\ge s\rangle,\qquad
 F_sN=\langle\tau+m:d(\tau)\le s\rangle.
\]
Both are $S_0$-filtrations, and their equal-jump strata are isomorphic:
\begin{equation}\label{eq:ordinary-strata}
 D_sT/D_{s+1}T\simeq F_sN/F_{s-1}N,
 \qquad\text{with degree shift }s.
\end{equation}
In fact, a surviving terminal action can only increase the jump, whereas
a surviving head action can only decrease it.  When the jumps agree, the
two actions survive together and have coefficient one.  To bound jumps,
write $Q=S_0/I_Q$ and $A=S_0/I_A$.  Since $Q_{\nu+1}=0$, $I_Q$, and hence
$N=I_Q/I_A$, is generated in degrees at most $\nu+1$.  A minimal head has
a positive-degree terminal, so its jump is at most $\nu$; head actions
propagate this bound.  A non-Ap\'ery head has degree at least three, which
also proves $A_2=Q_2$.

Let $p=\max\{u\ge0:\binom{u+1}{2}\le a\}$ and $r=\max(2,p)$.
We may replace $k$ by an infinite purely transcendental extension of the
same characteristic.  Green's degree-two restriction theorem
\cite{GreenRestriction} supplies independent linear forms
$E=(\lambda_1,\ldots,\lambda_r)$ such that $(A/EA)_{\ge2}=0$.
Indeed, if $a=\binom{p+1}{2}+b$, $0\le b\le p$, successive restrictions
bound the quadratic dimension after $u\le p$ cuts by
$\binom{p+1-u}{2}+\max(b-u,0)$, which vanishes at $u=p$.
If $p<2$, extend the cut to two forms.  Thus
\begin{equation}\label{eq:ordinary-cut-rank}
 2\le r\le e-1,\qquad
 r\ge3\Longrightarrow a_r:=\binom{r+1}{2}\le a\le\delta.
\end{equation}
Fix these forms once and put $P=k[\lambda_1,\ldots,\lambda_r]$.

For $2\le s\le\nu$, abbreviate $F_s=F_sN$ and define
\[
 U_s=A/F_{s-1},\quad X_s=(U_s)_{\ge s},\quad
 M_s=(U_s)_{\ge s+1},\quad K_s=F_s/F_{s-1}.
\]
Here $F_1=0$, $F_\nu=N$.  Jump-$s$ heads have degree at least $s+1$,
so these are actual exact sequences over the fixed ring $P$:
\begin{equation}\label{eq:ordinary-tower}
 0\longrightarrow K_s\longrightarrow M_s\longrightarrow X_{s+1}
 \longrightarrow0,\qquad X_{\nu+1}=0.
\end{equation}
The square-zero cut makes $U_s$ generated over $P$ in degrees zero and
one.  Its truncations $X_s,M_s$ are consequently generated only in
degrees $s,s+1$, respectively; their numbers of generators are $q_s$
and $\dim_k(X_s)_{s+1}$.
More explicitly, $(U_s)_n=E(U_s)_{n-1}$ for every $n\ge2$.
Iteration proves generation of each truncation in its bottom degree.
The degree-$s$ piece of $U_s$ agrees with $Q_s$, because every head
remaining outside $Q$ has jump at least $s$ and degree at least $s+1$.
When $\delta=0$, take $X_2=0$ in the formulas below.

\begin{lemma}\label{lem:ordinary-resources}
The finite sequence
$H_t=\sum_{s=2}^\nu\dim_k(X_s)_{s+t}$ satisfies, simultaneously,
\begin{equation}\label{eq:ordinary-resources}
 \begin{gathered}
 H_0=\delta,\qquad 0\le H_t\le B,\qquad
 S_H:=\sum_{t\ge0}H_t\le\delta\mathcal V,\\
 H_1=\delta-a+\ell\le B-a,\qquad H_1\le r\delta,
 \qquad \ell=\dim_kT_1.
 \end{gathered}
\end{equation}
For $\delta=0$ the tower and resource sequence are empty.
\end{lemma}

\begin{proof}
For a residue $\rho$ modulo $m$, let $b_\rho$ be the least sum of positive
offsets $g_i-m$ in that residue, and let $l_\rho$ be the shortest length
realizing $b_\rho$.  Write $\omega_\rho$ for the Ap\'ery representative,
$s_\rho=\operatorname{ord}(\omega_\rho)$, and
$L_\rho=(\omega_\rho-b_\rho)/m$.  Then
$1\le s_\rho\le L_\rho\le l_\rho$ for $\rho\ne0$.
Prefixes of a least-offset word are least-offset words and visit distinct
residues.  Along this path $L$ increases by at most one at each step,
and, after the first edge, cannot increase on arrival at an atom residue:
an increase there would decompose that atom. There are $e-1$ atom residues and $\delta$ other
nonzero residues.  Ordering nonzero residues by increasing $b_\rho$ gives
\[
 L_{\rho_j}-1\le\min(\delta,j-1),\qquad
 \sum_{\rho\ne0}(L_\rho-1)\le
 \sum_{j=1}^{m-1}\min(\delta,j-1)=\delta\mathcal V.
\]
Moreover, $v(q)=\operatorname{ord}(\omega_\rho+qm)$ eventually equals
$L_\rho+q$: the least-offset word gives this expression, and a longer
one would have smaller offset.  Telescoping therefore yields
\begin{equation}\label{eq:ordinary-terminal-budget}
 \sum_{\tau\equiv\rho\ (m)}(d(\tau)-1)=L_\rho-s_\rho.
\end{equation}
In $H$, an Ap\'ery coordinate contributes on $[0,s_\rho-2]$; a terminal
of degree $u$ and jump $d$ contributes its head on $[u,u+d-2]$.
These intervals are disjoint within each residue, since the next terminal
has degree at least $u+d$.  The zero residue contributes nothing.
Thus $H_t\le m-1$ and
$S_H=\sum_{\rho\ne0}(L_\rho-1)$.
At $t=0$ we count the $\delta$ Ap\'ery coordinates of degree at least two;
at $t=1$ we count $\delta-a$ higher Ap\'ery coordinates and $\ell$ linear
terminals.  Finally $\ell\le e-1$, and summing
$\dim_k(X_s)_{s+1}\le rq_s$ gives $H_1\le r\delta$.
\end{proof}

\subsection{From the tower to every relative row}

For a polynomial ring $\Lambda$, write
$b_i^\Lambda(Z)=\dim_k\operatorname{Tor}_i^\Lambda(Z,k)$ and
$c_j^\Lambda(Z)=\sum_{i=0}^j(-1)^{j-i}b_i^\Lambda(Z)$.
For a rank-zero module this is a syzygy rank and is nonnegative.

\begin{lemma}\label{lem:ordinary-harmonic}
If a finite-length $P$-module $Z$ is generated only in degree $b$ and
$h_t=\dim_kZ_{b+t}$, then
\[
 c_j^P(Z)\le\binom{r-1}{j}\sum_{t\ge0}\frac{j}{t+j}h_t
 \quad(1\le j<r),\qquad c_0^P(Z)=\mu_P(Z),\quad c_r^P(Z)=0.
\]
\end{lemma}

\begin{proof}
Shift to $b=0$ and take a homogeneous presentation
$P^q\twoheadrightarrow Z$. Choose a monomial order on its free module;
the monomials outside the initial module of the kernel give a basis of
$Z$, closed under divisors within each generator component. Label a
Koszul basis vector $m e_F$ by its augmented module monomial
$m\prod_{v\in F}x_v$. Reduction to standard monomials never increases
this label. Thus, in each internal degree, the actual differential is
block triangular in these labels, and its diagonal blocks are the
differential for the initial-module quotient. Choosing a maximal
nonsingular minor in each diagonal block and combining them gives a
block triangular minor of the actual matrix. Consequently its rank is
at least the sum of the diagonal ranks. The chain dimensions are
unchanged, so the identity for $c_j$ below reduces the desired upper
bound to standard monomials of monomial quotients.

In multidegree $\alpha$, $|\alpha|=t+j$, fix a variable $v$ in its support.
Every target row $(m,F)$ with $|F|=j$ and $v\notin F$ pairs with the
source column $(m/x_v,F\cup\{v\})$; divisors of standard monomials are
standard.  The selected minor is diagonal with entries $\pm1$.
Averaging these integer rank bounds with rational weights
$\alpha_v/(t+j)$ and summing gives
\[
 \operatorname{rank}\partial_{j+1}\big|_t\ge
 \binom{r-1}{j}\frac{t}{t+j}h_t.
\]
For $\lambda=\operatorname{length}Z$ and total boundary ranks $R_i$,
$c_j^P(Z)=\binom{r-1}{j}\lambda-R_{j+1}$, proving the claim.
The averaging uses rational numbers as numerical weights and hence
works in every characteristic.  The endpoint identities are the
generator count and the rank Euler identity.
\end{proof}

Set $R=k[x,E]$ and use $c_j(G)=c_j^R(G)$ from now on. The two-row $x$-Koszul
spectral sequence gives
$c_j(G)\le c_j^P(A)+c_{j-1}^P(T)$; its possible subsequent differential
subtracts a nonnegative rank.  For any exact sequence
$0\to K\to\mathcal M\to L\to0$, the Tor long exact sequence gives
\[
 c_j(\mathcal M)\le c_j(K)+c_j(L),\qquad
 c_j(K)\le c_j(\mathcal M)+c_{j+1}(L).
\]
For the first assertion, if $\rho_i$ denotes the rank of
$\operatorname{Tor}_i^P(A,k)\to\operatorname{Tor}_{i-2}^P(T,k)$, then
\[
 b_i^R(G)=b_i^P(A)+b_{i-1}^P(T)-\rho_i-\rho_{i+1},\qquad
 c_j(G)=c_j^P(A)+c_{j-1}^P(T)-\rho_{j+1}.
\]
For the second assertion, alternating the long exact sequence subtracts
either a connecting-map rank or its nonnegative complement in the
corresponding Tor group.  Thus all three comparisons preserve the
direction needed for an upper bound.
Apply these successively to $A$, the strata
\eqref{eq:ordinary-strata}, and \eqref{eq:ordinary-tower}.
Since $A/X_2\simeq P/P_+^2\oplus k(-1)^{e-r-1}$, we obtain
\begin{equation}\label{eq:ordinary-row}
 c_j(G)\le L_j+\sum_{t\ge0}W_j(t)H_t\quad(1\le j\le r),\qquad
 L_j=(e-1)\binom{r-1}{j}-\binom{r-1}{j+1},
\end{equation}
where
\[
 \begin{array}{ll}
 W_1(0)=r-1,&W_1(1)=(r+1)/2,\quad
 W_1(t)=(r-1)/(t+1)\quad(t\ge2),\\[2pt]
 W_j(0)=\binom{r-1}{j},&
 W_j(t)=\displaystyle\frac{j\binom{r-1}{j}}{t+j}
       +\frac{(j-1)\binom{r-1}{j-1}}{t+j-2}
       \quad(j\ge2,\ t\ge1).
 \end{array}
\]
Here the $X_s$ terms telescope, with last term $X_{\nu+1}=0$.
Lemma~\ref{lem:ordinary-harmonic} applies to $X_s$ and $M_s$;
for $j=1$ use the exact identity $\sum_s\mu_P(M_s)=H_1$.
The second denominator is shifted because $M_s$ starts one degree
later.  This also proves \eqref{eq:ordinary-row} at $j=r$.

\subsection{The target margin and the zeroth-homology correction}

Put $q=e-r$ and, with $e,r$ fixed, define
\[
 R_j(d)=(q+j)\binom{d+r-1}{j}
          +j\binom{d+r-1}{j+1},\qquad
 T_u(z)=1+\sum_{i\ge1}i\binom u{i+1}z^i.
\]
Pascal's identity gives
\[
 R_j(0)=L_j,\qquad
 R_j(d+1)-R_j(d)=(e+d)\binom{d+r-1}{j-1}.
\]
Consequently
\begin{equation}\label{eq:ordinary-margin}
 R_j(\delta)-L_j\ge
 \delta(\mathcal V+1)\binom{r-1}{j-1}.
\end{equation}
To explain why these are the required margins, extend $E$ to a basis
of the non-$x$ linear space, with $q-1$ remaining variables $Z$.
As a module over $R=k[x,E]$, $G$ has rank zero and $q$ generators.
Thus, writing $C(z)=\sum_{j=1}^rc_j(G)z^j$, its relative Betti polynomial
is $(1+z)(q+C(z))$.  Moreover
$G/(x,E)G\simeq k[Z]/(Z)^2$, whose Betti polynomial is $T_q$.
In the iterated Koszul complex, evaluate this zeroth homology row
exactly and bound the positive rows by their chain dimensions.  With
$\preceq$ denoting coefficientwise comparison, this gives
\begin{equation}\label{eq:ordinary-full-transfer}
 \mathcal B_G\preceq T_q+(1+z)^{q-1}
       \bigl((1+z)(q+C(z))-q\bigr)
       =T_{q+1}+(1+z)^qC(z).
\end{equation}
Here the relative degree-zero homology has dimension $q$, but the
remaining variables act on it: it is the square-zero quotient
$k[Z]/(Z)^2$, not a trivial $q$-dimensional coefficient space. Its
contribution must therefore be $T_q$, rather than
$q(1+z)^{q-1}$. Subtracting $q$ inside the parentheses removes the
degree-zero term before the positive relative rows are bounded by
chain dimensions. This prevents double counting and gives the exact
zeroth row $1$. No spectral-sequence degeneration is used.
Summing $T_{u+1}-T_u=uz(1+z)^{u-1}$ yields
\begin{equation}\label{eq:ordinary-target-difference}
 T_m-T_{q+1}=(1+z)^q
       \sum_{j=1}^{\delta+r-1}R_j(\delta)z^j.
\end{equation}
It therefore suffices to prove $c_j(G)\le R_j(\delta)$ for $1\le j\le r$;
the additional target coefficients are nonnegative.
Indeed the coefficient before factoring out $(1+z)^q$ is obtained from
$\sum_{u=0}^{\delta+r-2}(q+1+u)\binom{u}{j-1}$.
The hockey-stick identity and
$u\binom{u}{j-1}=(j-1)\binom{u}{j-1}+j\binom{u}{j}$ evaluate it as
$R_j(\delta)$.  This calculation includes $\delta=0$: out-of-range
binomial coefficients vanish, and the relative top margin is zero.

\subsection{The numerical inequalities}

For $r=2$, \eqref{eq:ordinary-row} and
$H_1\le\min(S_H-\delta,2\delta)$ give
\[
 c_1-L_1\le\min\{\tfrac32S_H-\tfrac12\delta,
                         \tfrac52\delta+\tfrac12S_H\}
 \le\delta(\mathcal V+1).
\]
Use the first estimate for $\mathcal V\le3$ and the second for
$\mathcal V\ge3$.  The top row satisfies $c_2\le S_H-\delta$.
For $r=3$, $e\ge4$, $\delta\ge6$, and $\mathcal V\ge11/2$.
The first row is at most
$L_1+\frac43\delta+\frac43H_1+\frac23S_H$.
Using $H_1\le B-6$, three times its margin is bounded below by
\[
 \delta(\mathcal V-1)-4(B-6)
 =zs+\tfrac12z^2+\tfrac72z+2s+15>0,
 \quad z=\delta-6,\ s=e-4.
\]
The second row is at most
$L_2+\frac32S_H-\frac12\delta+\frac76H_1
 \le L_2+\frac32S_H+3\delta
 \le L_2+2\delta(\mathcal V+1)$;
the top row is again at most $S_H-\delta$.
Thus \eqref{eq:ordinary-margin} proves all low-rank comparisons.

For $r\ge4$ define the single scalar
\[
 F_r=(r-1)\delta+\frac{r+1}{2}H_1
              +\sum_{t\ge2}\frac{r-1}{t+1}H_t.
\]
Dividing $W_j(t)$ by $\binom{r-1}{j-1}$ bounds it by the corresponding
coefficient of $F_r$.  At $t\ge2$ the normalized weight is
$(r-j)/(t+j)+(j-1)/(t+j-2)$.
For $j\ge3$ both denominators are at least $t+1$; for $j=2$ its difference
from $(r-1)/(t+1)$ is
$[2-(r-3)t]/[t(t+1)(t+2)]\le0$.
The cases $t=0,1$ follow directly from the binomial identities.
These estimates depend only on the displayed resource inequalities.
Hence all rows, including the top, follow once
\begin{equation}\label{eq:ordinary-scalar}
 F_r\le\delta(\mathcal V+1).
\end{equation}

Here is a common cutoff argument for this last inequality.
Set $a_*=a_r$ and choose an integer $n\ge r$; write
$\alpha=(r-1)/(n+1)$, $\beta=1-\alpha$,
$c=(r+1)/2-\alpha$, and $\mathsf h_n=\sum_{j=1}^n1/j$.
Subtracting $\alpha$ from every weight and using
\eqref{eq:ordinary-resources} gives
\begin{equation}\label{eq:ordinary-cutoff}
 \begin{split}
 F_r\le U(e,\delta):={}&\alpha\delta\mathcal V
 +(r-1-\alpha)\delta+c(B-a_*)\\
 &+B\sum_{j=2}^{n-1}\left(\frac{r-1}{j+1}-\alpha\right).
 \end{split}
\end{equation}
The omitted residual weights for $j\ge n$ are nonpositive; no support
bound is assumed.  Define $\Delta(e,d)=d(e+d/2-1/2)-U(e,d)$.
For any base $e_0$, put $g=\Delta(e_0,a_*)$,
$A_*=\Delta(e_0+1,a_*)-g$, and $b=\Delta(e_0,a_*+1)-g$.
Summing the finite harmonic expression gives
\begin{equation}\label{eq:ordinary-gap-shift}
 \begin{gathered}
 A_*=a_*-1-(r-1)\left(\mathsf h_n-2+\frac{a_*+2}{n+1}\right),\\
 \Delta(e_0+s,a_*+t)
 =g+sA_*+tb+\beta\left(\frac{t(t-1)}2+st\right).
 \end{gathered}
\end{equation}
This is an identity in rational $s,t$.

For $4\le r\le8$ use $e_0=r+1$ and the following exact values:
\[
 \begin{array}{c|r|r|r|r|r|r}
 r&a_*&n&L&Lg&LA_*&Lb\\\hline
 4&10&7&280&2443&762&1007\\
 5&15&10&6930&39380&28429&32209\\
 6&21&13&72072&-7479&424147&439591\\
 7&28&17&2042040&-12196485&17076937&17076937\\
 8&36&21&739024&-5070732&8418651&8250691
 \end{array}
\]
All $A_*,b,\beta$ are positive.  The first two seeds are positive;
in the last three rows $g+A_*,g+b>0$.  Thus
\eqref{eq:ordinary-gap-shift} is positive for nonnegative integer $s,t$,
except possibly $s=t=0$ in the last three rows.
Those corners cannot occur.  Indeed, if $r=e-1\ge3$, then
$a=\binom e2$, so all quadratic products of the non-$m$ atoms give
distinct nonzero degree-two Ap\'ery coordinates.  The residues of
$\{0,g_2,\ldots,g_e\}$ consequently form a Sidon set modulo $m$.
Its $e(e-1)$ ordered nonzero differences are distinct, giving
$m\ge e(e-1)+1$ and $\delta\ge(e-1)^2>a_*$.
This proves \eqref{eq:ordinary-scalar} in these ranks.

For completeness, the unbounded range has a uniform argument.
Let $r\ge9$, $a_*=r(r+1)/2$, and $D=\lfloor a_*/2\rfloor+1$;
use $n=D$ in \eqref{eq:ordinary-cutoff}.  Then
$a_*+1\le2D\le a_*+2$ and $D\ge r$.
For $\mathsf b(r)=r/4+1+9/(2r)$,
\[
 \mathsf h_{D_{r+1}}-\mathsf h_{D_r}
 \le\frac{r+2}{a_r+3}<\mathsf b(r+1)-\mathsf b(r)\quad(r\ge10).
\]
The strict inequality follows on putting $x=r-10$ from the positive
numerator $x^4+34x^3+385x^2+1472x+112$ over
$4r(r+1)(r^2+r+6)$.
Together with
$\mathsf b(10)-\mathsf h_{28}=1833472237/80313433200>0$, this proves
$\mathsf h_D<\mathsf b(r)$ for every $r\ge10$.
Define
\[
 T_r=\frac{a_*(a_*+22-r)}{2(r-1)(a_*+9)}
       +\frac{2r-3}{r-1}
       -\frac{a_*(a_*+21)+36}{2(D+1)(a_*+9)}.
\]
Direct summation in \eqref{eq:ordinary-cutoff} gives
$g_0:=\Delta(10,a_*)=(r-1)(a_*+9)(T_r-\mathsf h_D)$.
Replace the last denominator of $T_r$ by $(a_*+3)(a_*+9)$ to obtain
$T_-\le T_r$.  Its difference from $\mathsf b(r)$ is
\[
 T_--\mathsf b(r)=
 \frac{r^3[(r-10)(2r+15)+2]+r[111(r-10)+210]+972}
 {2r(r-1)(r^2+r+6)(r^2+r+18)}>0.
\]
Thus $g_0>0$ for $r\ge10$; for $r=9$, direct substitution gives
$g_0=5845533/7436429>0$.

It remains to extend this seed in both parameters.
With $A_*=\Delta(11,a_*)-g_0$, $b_0=\Delta(10,a_*+1)-g_0$, put
\[
 C_*=\frac{a_*[-a_*(D-r+2)-D(r-4)+3]}{2(D+1)}.
\]
The identities $g_0=(a_*+9)A_*+C_*$ and
$b_0=A_*+11-r-8\alpha$ give $A_*>-C_* /(a_*+9)>0$.
For $F(D)=-C_* /(a_*+9)+11-r-8\alpha$, set
\[
 F_{\min}=\frac{a_*}2-\frac{3r}2+\frac{11}2
       +\frac{3(5r+31)}{2(a_*+9)}-\frac{17(r-1)}{a_*+3}.
\]
The floor comparison is the identity
\[
 F(D)-F_{\min}=
 \frac{(2D-a_*-1)(r-1)(a_*^2+17a_*+144)}
 {2(a_*+9)(D+1)(a_*+3)}\ge0.
\]
Since $a_*\ge4r$, $F_{\min}>11/2-17/4=5/4$;
hence $b_0>5/4$.  Finally $0<\beta<1$, and
\eqref{eq:ordinary-gap-shift}, recentered at $e_0=10$, reads
\[
 \Delta(10+s,a_*+t)=g_0+sA_*+t(b_0-\beta/2)
                       +\beta(t^2/2+st)>0
 \qquad(s,t\ge0).
\]
This proves the scalar inequality for all rational $e\ge10$,
$\delta\ge a_*$, and therefore for every actual cut with $r\ge9$.

\begin{proof}[Proof of Proposition~\ref{prop:ordinary}]
The low-rank estimates and the scalar inequality, combined with
\eqref{eq:ordinary-row} and \eqref{eq:ordinary-margin}, give
$c_j(G)\le R_j(\delta)$ for every $1\le j\le r$.
Equations~\eqref{eq:ordinary-full-transfer}--\eqref{eq:ordinary-target-difference}
yield $\beta_i(G;k)\le i\binom m{i+1}$.
The width assumption identifies the minimal generators of the interval
completion as $m,m+1,\ldots,2m-1$; its Betti numbers are those in
\eqref{eq:ordinary-capacity}, by \cite{APFormula,APCone}.
Finally, scalar extension preserves the original monomial coordinates
and the ranks in a minimal graded free resolution.  Thus the bound
proved after the purely transcendental extension descends to the
original field.  The algebraic and completed tangent cones have the
same graded coordinates by Proposition~\ref{prop:bridge}; the conclusion
therefore concerns the stated completed semigroup rings and their own
minimal presentations.
\end{proof}
